% Generated by tools/simcode_parameters.py from FML/SimulationCode/PARAMETERS.md. Do not edit.
\chapter{SimulationCode parameters}\label{app:simcode-parameters}

Every parameter of SimulationCode (\cref{ch:simulation-code}), as the code defines it in \file{src/ParameterSchema.h} and \file{SimulationParameters.h}. This appendix is generated from \file{SimulationCode/PARAMETERS.md}, which \code{./nbody --parameters PARAMETERS.md} writes, so it lists exactly the parameters, defaults, conditions and allowed values of the code. Each entry gives the type and the default value (\emph{required}: no default), when the parameter is read, and its meaning.

\section{General}

{\small\begin{longtable}{@{}>{\raggedright\arraybackslash}p{4.3cm}>{\raggedright\arraybackslash}p{2.8cm}>{\raggedright\arraybackslash}p{8.1cm}@{}}
\toprule Parameter (type, default) & Used when & Description \\ \midrule
\endfirsthead
\toprule Parameter (type, default) & Used when & Description \\ \midrule
\endhead
\bottomrule
\endlastfoot
\param{all_parameters_must_be_in_file}\newline{\footnotesize bool, \code{false}} & always & If true every parameter that is used must be in the parameterfile (no default values are used) \\[2pt]
\end{longtable}}

\section{Simulation}

{\small\begin{longtable}{@{}>{\raggedright\arraybackslash}p{4.3cm}>{\raggedright\arraybackslash}p{2.8cm}>{\raggedright\arraybackslash}p{8.1cm}@{}}
\toprule Parameter (type, default) & Used when & Description \\ \midrule
\endfirsthead
\toprule Parameter (type, default) & Used when & Description \\ \midrule
\endhead
\bottomrule
\endlastfoot
\param{simulation_boxsize}\newline{\footnotesize double, \emph{required}} & always & Boxsize of the simulation (Mpc/h) \\[2pt]
\param{simulation_name}\newline{\footnotesize string, \emph{required}} & always & Label used in the output filenames \\[2pt]
\param{simulation_use_cola}\newline{\footnotesize bool, \code{false}} & always & Use the COLA method \\[2pt]
\param{simulation_use_scaledependent_cola}\newline{\footnotesize bool, \code{true}} & always & If the growth is scale-dependent (massive neutrinos, k-dependent gravity models): use scale-dependent COLA. If false we use the k = 0 limit of the growth factors when doing COLA. With scale-independent growth (or without COLA) it is set to false (standard COLA gives the same result at a much lower cost) \\[2pt]
\param{simulation_enforce_LPT_trajectories}\newline{\footnotesize bool, \code{false}} & \code{simulation_use_cola} = true & Do not compute the PM force, just let particles follow the LPT trajectories \\[2pt]
\param{simulation_cola_LPT_order}\newline{\footnotesize int, \code{2}} & \code{simulation_use_cola} = true & The LPT order of COLA (1 or 2). Has to be <= \code{ic_LPT_order.} If < 1 COLA is turned off \\[2pt]
\param{simulation_fftw_planner}\newline{\footnotesize string, \code{"estimate"}} & always & How FFTW makes its plans. estimate: no measuring (fast to start, slower transforms). measure, patient, exhaustive: time candidate algorithms the first time each grid size is transformed (faster transforms, often 1.5-2x for large grids, but planning can take seconds to minutes). The plans made are written to \code{fftw_wisdom_[simulation_name].txt} in the output folder (not for estimate): give that file as \code{simulation_fftw_wisdom_file} to skip the planning in later runs. Options: \code{estimate}, \code{measure}, \code{patient}, \code{exhaustive} \\[2pt]
\param{simulation_fftw_wisdom_file}\newline{\footnotesize string, \code{""}} & always & FFTW wisdom (plans) to read at the start, e.g. the \code{fftw_wisdom_[simulation_name].txt} written by an earlier run ("": none). Plans in it are used for the transforms they match (same FFTW version, precision, grid size and number of threads; a different machine or thread count gives plans that are valid but possibly slow). Other transforms are planned with \code{simulation_fftw_planner} \\[2pt]
\end{longtable}}

\section{Cosmology}

{\small\begin{longtable}{@{}>{\raggedright\arraybackslash}p{4.3cm}>{\raggedright\arraybackslash}p{2.8cm}>{\raggedright\arraybackslash}p{8.1cm}@{}}
\toprule Parameter (type, default) & Used when & Description \\ \midrule
\endfirsthead
\toprule Parameter (type, default) & Used when & Description \\ \midrule
\endhead
\bottomrule
\endlastfoot
\param{cosmology_model}\newline{\footnotesize string, \code{"LCDM"}} & always & The cosmology (background evolution). Options: \code{LCDM}, \code{w0waCDM}, \code{DGP}, \code{JBD}, \code{Quintessence}, \code{EDE}, \code{DarkScattering}, \code{CovariantDarkScattering}, \code{CoupledQuintessence}, \code{FadingDarkSector}, \code{IDE}, \code{CubicGalileon}, \code{Horndeski}, \code{Tabulated} \\[2pt]
\param{cosmology_Omegab}\newline{\footnotesize double, \emph{required}} & always & Baryon density parameter \\[2pt]
\param{cosmology_OmegaMNu}\newline{\footnotesize double, \code{0}} & always & Massive neutrino density parameter (3 equal masses; or use \code{cosmology_neutrino_masses_eV}) \\[2pt]
\param{cosmology_neutrino_masses_eV}\newline{\footnotesize double list, \texttt{\{\}}} & always & The masses of the (up to 3) neutrinos in eV, e.g. \{0.0, 0.0086, 0.05\} (0 or not listed: massless). Instead of \code{cosmology_OmegaMNu} (which then follows from the masses) \\[2pt]
\param{cosmology_OmegaCDM}\newline{\footnotesize double, \emph{required}} & always & CDM density parameter \\[2pt]
\param{cosmology_OmegaK}\newline{\footnotesize double, \emph{required}} & always & Curvature density parameter (-k/H0\textasciicircum{}2) \\[2pt]
\param{cosmology_Neffective}\newline{\footnotesize double, \emph{required}} & always & Effective number of relativistic species \\[2pt]
\param{cosmology_TCMB_kelvin}\newline{\footnotesize double, \emph{required}} & always & Temperature of the CMB today (Kelvin) \\[2pt]
\param{cosmology_h}\newline{\footnotesize double, \emph{required}} & always & Hubble parameter H0 = 100h km/s/Mpc \\[2pt]
\param{cosmology_JBD_wBD}\newline{\footnotesize double, \emph{required}} & \code{cosmology_model} = JBD & The JBD parameter (wBD -> infinity gives GR) \\[2pt]
\param{cosmology_JBD_GeffG_today}\newline{\footnotesize double, \code{1}} & \code{cosmology_model} = JBD & The value of G today relative to Newtons constant (\code{GeffG_today} != 1 effectively gives a different Newtons constant) \\[2pt]
\param{cosmology_JBD_density_parameter_definition}\newline{\footnotesize string, \code{"hi-class"}} & \code{cosmology_model} = JBD & Definition of the density parameters Omega = 8pi \code{G_}\emph{ rho / 3H0\textasciicircum{}2 (which G\_} to use). Options: \code{Gbare}, \code{Gnewton}, \code{Gtoday}, \code{hi-class} \\[2pt]
\param{cosmology_w0}\newline{\footnotesize double, \emph{required}} & \code{cosmology_model} = w0waCDM & Dark energy equation of state w(a) = w0 + wa(1-a): w0 \\[2pt]
\param{cosmology_wa}\newline{\footnotesize double, \emph{required}} & \code{cosmology_model} = w0waCDM & Dark energy equation of state w(a) = w0 + wa(1-a): wa \\[2pt]
\param{cosmology_quintessence_potential}\newline{\footnotesize string, \emph{required}} & \code{cosmology_model} = Quintessence & The potential V = V0 u(phi), phi in units of Mpl: exponential u = exp(-lambda phi), \code{inverse_power_law} u = phi\textasciicircum{}-alpha, pngb u = 1 + cos(phi/f), hilltop u = (1 - (phi/mu)\textasciicircum{}p)\textasciicircum{}2. V0 is found so that OmegaDE today = 1 - OmegaM - OmegaR - OmegaK. Options: \code{exponential}, \code{inverse_power_law}, \code{pngb}, \code{hilltop} \\[2pt]
\param{cosmology_quintessence_phi_ini}\newline{\footnotesize double, \code{0}} & \code{cosmology_model} = Quintessence & The initial value of the field (units of Mpl; deep in the radiation era) \\[2pt]
\param{cosmology_quintessence_dphidlna_ini}\newline{\footnotesize double, \code{0}} & \code{cosmology_model} = Quintessence & The initial dphi/dlna (0: frozen) \\[2pt]
\param{cosmology_quintessence_lambda}\newline{\footnotesize double, \emph{required}} & \code{cosmology_model} = Quintessence and \code{cosmology_quintessence_potential} = exponential & The slope lambda \\[2pt]
\param{cosmology_quintessence_alpha}\newline{\footnotesize double, \emph{required}} & \code{cosmology_model} = Quintessence and \code{cosmology_quintessence_potential} = \code{inverse_power_law} & The power alpha \\[2pt]
\param{cosmology_quintessence_f}\newline{\footnotesize double, \emph{required}} & \code{cosmology_model} = Quintessence and \code{cosmology_quintessence_potential} = pngb & The scale f (units of Mpl) \\[2pt]
\param{cosmology_quintessence_mu}\newline{\footnotesize double, \emph{required}} & \code{cosmology_model} = Quintessence and \code{cosmology_quintessence_potential} = hilltop & The position of the minimum mu (units of Mpl) \\[2pt]
\param{cosmology_quintessence_p}\newline{\footnotesize double, \emph{required}} & \code{cosmology_model} = Quintessence and \code{cosmology_quintessence_potential} = hilltop & The power p \\[2pt]
\param{cosmology_ede_fEDE}\newline{\footnotesize double, \emph{required}} & \code{cosmology_model} = EDE & Early dark energy: the maximum of \code{rho_EDE} / \code{rho_total} \\[2pt]
\param{cosmology_ede_zc}\newline{\footnotesize double, \emph{required}} & \code{cosmology_model} = EDE & Early dark energy: the redshift of the maximum \\[2pt]
\param{cosmology_ede_theta_ini}\newline{\footnotesize double, \code{2.8}} & \code{cosmology_model} = EDE & Early dark energy: the initial angle \code{phi_ini} / f \\[2pt]
\param{cosmology_ede_n}\newline{\footnotesize double, \code{3}} & \code{cosmology_model} = EDE & Early dark energy: V = m\textasciicircum{}2 f\textasciicircum{}2 (1 - cos(phi/f))\textasciicircum{}n \\[2pt]
\param{cosmology_tabulated_filename}\newline{\footnotesize string, \emph{required}} & \code{cosmology_model} = Tabulated & The file with the background \\[2pt]
\param{cosmology_tabulated_format}\newline{\footnotesize string, \emph{required}} & \code{cosmology_model} = Tabulated & H: columns z, H/H0. CLASS: a CLASS background file. w: columns a, w(a) of the dark energy. Options: \code{H}, \code{CLASS}, \code{w} \\[2pt]
\param{cosmology_ds_w}\newline{\footnotesize double, \emph{required}} & \code{cosmology_model} in \{DarkScattering, CovariantDarkScattering\} & Dark scattering: the (constant) equation of state w of the dark energy (!= -1) \\[2pt]
\param{cosmology_ds_cs2}\newline{\footnotesize double, \code{1}} & \code{cosmology_model} in \{DarkScattering, CovariantDarkScattering\} & Dark scattering: the rest-frame sound speed squared of the dark energy (only for the IC from the internal Einstein-Boltzmann solver) \\[2pt]
\param{cosmology_cds_alpha}\newline{\footnotesize double, \emph{required}} & \code{cosmology_model} = CovariantDarkScattering & Covariantised dark Thomson-like scattering: alpha = 8 pi G \code{alpha_bar} / (3 H0\textasciicircum{}3) of the interaction Q\textasciicircum{}mu = \code{alpha_bar} (\code{u_DE}\textasciicircum{}mu - \code{u_c}\textasciicircum{}mu) (0: wCDM). The dark matter feels the drag alpha a\textasciicircum{}3 / (\code{Omega_cdm} E(a)) per ln a (the CDM+baryon particles \code{f_c} = OmegaCDM / OmegaCB of it) \\[2pt]
\param{cosmology_ds_xi}\newline{\footnotesize double, \emph{required}} & \code{cosmology_model} = DarkScattering & Dark scattering: xi = \code{sigma_D} / \code{m_c} in barn/GeV (0: wCDM). The dark matter feels the drag (1 + w) c xi \code{rho_DE} / H per ln a (the CDM+baryon particles \code{f_c} = OmegaCDM / OmegaCB of it) \\[2pt]
\param{cosmology_cq_lambda}\newline{\footnotesize double, \emph{required}} & \code{cosmology_model} = CoupledQuintessence & Coupled quintessence: the slope lambda of the potential V \textasciitilde{} exp(-lambda phi) (phi in reduced Planck masses) \\[2pt]
\param{cosmology_cq_beta}\newline{\footnotesize double, \emph{required}} & \code{cosmology_model} = CoupledQuintessence & Coupled quintessence: the coupling beta (Amendola's sqrt(2/3) beta): the CDM particles have the mass m \textasciitilde{} exp(-beta phi) and feel G (1 + 2 beta\textasciicircum{}2) from the CDM. 0: uncoupled \\[2pt]
\param{cosmology_fds_c}\newline{\footnotesize double, \emph{required}} & \code{cosmology_model} = FadingDarkSector & Fading Dark Sector (arXiv:2611.17422): c of \code{V_DE} = V0 exp(-c phi) (phi in reduced Planck masses; = \code{cosmology_cq_lambda}). The paper's best fit: -0.85 \\[2pt]
\param{cosmology_fds_cprime}\newline{\footnotesize double, \emph{required}} & \code{cosmology_model} = FadingDarkSector & Fading Dark Sector: c' of \code{m_DM} = m0 exp(-c' phi) (= \code{cosmology_cq_beta}): the CDM feels G (1 + 2 c'\textasciicircum{}2) from the CDM. The paper's best fit: 0.05 \\[2pt]
\param{cosmology_horndeski_filename}\newline{\footnotesize string, \emph{required}} & \code{cosmology_model} = Horndeski & Generic reduced Horndeski: the \code{_gravity.txt} table of \code{scripts/horndeski_frontend.py} (its H(a), and the alphas for the internal Einstein-Boltzmann solver); use it with \code{gravity_model} = Horndeski \\[2pt]
\param{cosmology_ide_w}\newline{\footnotesize double, \emph{required}} & \code{cosmology_model} = IDE & Interacting dark energy: the (constant) equation of state w of the dark energy (!= -1) \\[2pt]
\param{cosmology_ide_xi}\newline{\footnotesize double, \emph{required}} & \code{cosmology_model} = IDE & Interacting dark energy: the coupling xi of the energy exchange Q = xi H \code{rho_DE} along the dark-matter 4-velocity (> 0: energy flows to the dark matter; 0: wCDM). Needs the doom factor xi / (3 (1 + w)) <= 0 (xi < 0 for w > -1, xi > 0 for w < -1) \\[2pt]
\param{cosmology_ide_cs2}\newline{\footnotesize double, \code{1}} & \code{cosmology_model} = IDE & Interacting dark energy: the rest-frame sound speed squared of the dark energy (only for the IC from the internal Einstein-Boltzmann solver) \\[2pt]
\param{cosmology_dgp_OmegaRC}\newline{\footnotesize double, \emph{required}} & \code{cosmology_model} = DGP & The cross-over scale of the self-accelerating DGP model (\code{Omega_rc} = 1/(4 rc\textasciicircum{}2 H0\textasciicircum{}2)). A cosmological constant makes up E(0) = 1 (a warning if it is not zero; no dark energy for \code{Omega_rc} = (1 - \code{Omega_matter} - \code{Omega_radiation})\textasciicircum{}2/4) \\[2pt]
\end{longtable}}

\section{Gravity model}

{\small\begin{longtable}{@{}>{\raggedright\arraybackslash}p{4.3cm}>{\raggedright\arraybackslash}p{2.8cm}>{\raggedright\arraybackslash}p{8.1cm}@{}}
\toprule Parameter (type, default) & Used when & Description \\ \midrule
\endfirsthead
\toprule Parameter (type, default) & Used when & Description \\ \midrule
\endhead
\bottomrule
\endlastfoot
\param{gravity_model}\newline{\footnotesize string, \code{"GR"}} & always & The gravity model (growth of perturbations and forces). Options: \code{GR}, \code{f(R)}, \code{DGP}, \code{CubicGalileon}, \code{Horndeski}, \code{JBD}, \code{Symmetron}, \code{Dilaton}, \code{mbeta}, \code{mu}, \code{CoupledQuintessence}, \code{FadingDarkSector} \\[2pt]
\param{gravity_model_mu_parametrization}\newline{\footnotesize string, \emph{required}} & \code{gravity_model} = mu & Poisson equation with \code{G_eff/G} = mu(a,k), no screening. DE: mu = 1 + mu0 \code{Omega_DE}(a)\code{/Omega_DE}(a=1). BZ: mu = (1 + beta1 lambda1\textasciicircum{}2 k\textasciicircum{}2 a\textasciicircum{}s) / (1 + lambda1\textasciicircum{}2 k\textasciicircum{}2 a\textasciicircum{}s) (Bertschinger \& Zukin 2008). table: mu(a) or mu(a,k) from a file. Options: \code{DE}, \code{BZ}, \code{table} \\[2pt]
\param{gravity_model_mu_mu0}\newline{\footnotesize double, \emph{required}} & \code{gravity_model} = mu and \code{gravity_model_mu_parametrization} = DE & mu0 in mu = 1 + mu0 \code{Omega_DE}(a)\code{/Omega_DE}(a=1) \\[2pt]
\param{gravity_model_mu_BZ_params}\newline{\footnotesize double list, \emph{required}} & \code{gravity_model} = mu and \code{gravity_model_mu_parametrization} = BZ & \{beta1, lambda1 (Mpc/h), s\} in mu = (1 + beta1 lambda1\textasciicircum{}2 k\textasciicircum{}2 a\textasciicircum{}s) / (1 + lambda1\textasciicircum{}2 k\textasciicircum{}2 a\textasciicircum{}s) (k in h/Mpc) \\[2pt]
\param{gravity_model_mu_filename}\newline{\footnotesize string, \emph{required}} & \code{gravity_model} = mu and \code{gravity_model_mu_parametrization} = table & File with rows 'a mu' or 'a k mu' (k in h/Mpc, a regular grid in a and k; e.g. from \code{hi_class}). Splined in log a, log k; the closest value outside \\[2pt]
\param{gravity_model_mbeta_params}\newline{\footnotesize double list, \emph{required}} & \code{gravity_model} = mbeta & The parameters (beta0, n, m0, m) of beta(a) = beta0 a\textasciicircum{}n and m(a) = m0 H0 a\textasciicircum{}m \\[2pt]
\param{gravity_model_fofr_fofr0}\newline{\footnotesize double, \emph{required}} & \code{gravity_model} = f(R) & The value of \code{f_R} today \\[2pt]
\param{gravity_model_fofr_nfofr}\newline{\footnotesize double, \code{1}} & \code{gravity_model} = f(R) & The index n of the Hu-Sawicky model \\[2pt]
\param{gravity_model_symmetron_assb}\newline{\footnotesize double, \code{0.5}} & \code{gravity_model} = Symmetron & Symmetry breaking scale factor (no fifth-force for a < assb) \\[2pt]
\param{gravity_model_symmetron_beta}\newline{\footnotesize double, \code{1}} & \code{gravity_model} = Symmetron & Coupling strength relative to gravity \\[2pt]
\param{gravity_model_symmetron_L_mpch}\newline{\footnotesize double, \code{1}} & \code{gravity_model} = Symmetron & Range of the fifth force in vacuum (Mpc/h; for (N, M) = (2, 4) also the range in the background today if \code{a_ssb}\textasciicircum{}3 << 1) \\[2pt]
\param{gravity_model_symmetron_N}\newline{\footnotesize double, \code{2}} & \code{gravity_model} = Symmetron & The power N in V = -(\code{phi/phi_}\emph{)\textasciicircum{}N / N + (phi/phi\_})\textasciicircum{}M / M and beta(phi) \textasciitilde{} phi\textasciicircum{}(N-1) (even integer >= 2; Brax et al. 2012, arXiv:1206.3568). (N, M) = (2, 4) is the original symmetron \\[2pt]
\param{gravity_model_symmetron_M}\newline{\footnotesize double, \code{4}} & \code{gravity_model} = Symmetron & The power M in the potential (even integer > N) \\[2pt]
\param{gravity_model_dilaton_A2}\newline{\footnotesize double, \code{500000}} & \code{gravity_model} = Dilaton & A2 in the coupling A(phi) = 1 + A2 phi\textasciicircum{}2 / (2 Mpl\textasciicircum{}2) (A2 >> 1; Brax et al. 2012, arXiv:1206.3568). The screening threshold is \code{Phi_crit} = 1 / (2 A2) \\[2pt]
\param{gravity_model_dilaton_beta0}\newline{\footnotesize double, \code{0.5}} & \code{gravity_model} = Dilaton & The coupling today beta0 in beta(a) = beta0 exp[s (a\textasciicircum{}(2r-3) - 1) / (2r - 3)], s = 9 OmegaM A2 xi\textasciicircum{}2 \\[2pt]
\param{gravity_model_dilaton_r}\newline{\footnotesize double, \code{1}} & \code{gravity_model} = Dilaton & The index r in the mass m(a) = m0 a\textasciicircum{}-r (r <= 3/2; 3/2 is the original dilaton with beta = beta0 a\textasciicircum{}s) \\[2pt]
\param{gravity_model_dilaton_xi}\newline{\footnotesize double, \code{0.001}} & \code{gravity_model} = Dilaton & xi = H0 / m0: the range of the fifth force today is 2998 xi Mpc/h \\[2pt]
\param{gravity_model_dgp_rcH0overc}\newline{\footnotesize double, \emph{required}} & \code{gravity_model} = DGP & The cross-over scale rc H0 / c \\[2pt]
\param{gravity_model_horndeski_filename}\newline{\footnotesize string, \code{""}} & \code{gravity_model} = Horndeski & Generic reduced Horndeski gravity: the table made by \code{scripts/horndeski_frontend.py} from the Lagrangian (K, G3, G4); empty: that of \code{cosmology_model} = Horndeski (\code{cosmology_horndeski_filename}). With \code{cosmology_model} = Tabulated (the frontend's \code{_background.txt}, format H) it must be given \\[2pt]
\param{gravity_model_screening}\newline{\footnotesize bool, \code{true}} & \code{gravity_model} in \{mbeta, f(R), Symmetron, Dilaton, DGP, CubicGalileon, Horndeski\} & Use the approximate screening model (otherwise linear) \\[2pt]
\param{gravity_model_dgp_smoothing_filter}\newline{\footnotesize string, \code{"tophat"}} & \code{gravity_model} in \{mbeta, f(R), Symmetron, Dilaton, DGP, CubicGalileon, Horndeski\} and \code{gravity_model_screening} = true and \code{gravity_model} = DGP & Screening approximation: the filter used to smooth the density field. Options: \code{tophat}, \code{sharpk}, \code{gaussian} \\[2pt]
\param{gravity_model_dgp_smoothing_scale_over_boxsize}\newline{\footnotesize double, \code{0}} & \code{gravity_model} in \{mbeta, f(R), Symmetron, Dilaton, DGP, CubicGalileon, Horndeski\} and \code{gravity_model_screening} = true and \code{gravity_model} = DGP & Screening approximation: the smoothing scale R/boxsize (0: no smoothing; 1 smooths over the whole box, which makes the screening very weak) \\[2pt]
\param{gravity_model_cubic_galileon_smoothing_filter}\newline{\footnotesize string, \code{"tophat"}} & \code{gravity_model} in \{mbeta, f(R), Symmetron, Dilaton, DGP, CubicGalileon, Horndeski\} and \code{gravity_model_screening} = true and \code{gravity_model} = CubicGalileon & Screening approximation: the filter used to smooth the density field. Options: \code{tophat}, \code{sharpk}, \code{gaussian} \\[2pt]
\param{gravity_model_cubic_galileon_smoothing_scale_over_boxsize}\newline{\footnotesize double, \code{0}} & \code{gravity_model} in \{mbeta, f(R), Symmetron, Dilaton, DGP, CubicGalileon, Horndeski\} and \code{gravity_model_screening} = true and \code{gravity_model} = CubicGalileon & Screening approximation: the smoothing scale R/boxsize (0: no smoothing) \\[2pt]
\param{gravity_model_horndeski_smoothing_filter}\newline{\footnotesize string, \code{"gaussian"}} & \code{gravity_model} in \{mbeta, f(R), Symmetron, Dilaton, DGP, CubicGalileon, Horndeski\} and \code{gravity_model_screening} = true and \code{gravity_model} = Horndeski & Screening approximation: the filter used to smooth the density field. Options: \code{tophat}, \code{sharpk}, \code{gaussian} \\[2pt]
\param{gravity_model_horndeski_smoothing_scale_over_boxsize}\newline{\footnotesize double, \code{0}} & \code{gravity_model} in \{mbeta, f(R), Symmetron, Dilaton, DGP, CubicGalileon, Horndeski\} and \code{gravity_model_screening} = true and \code{gravity_model} = Horndeski & Screening approximation: the smoothing scale R/boxsize (0: no smoothing) \\[2pt]
\param{gravity_model_screening_enforce_largescale_linear}\newline{\footnotesize bool, \code{false}} & \code{gravity_model} in \{mbeta, f(R), Symmetron, Dilaton, DGP, CubicGalileon, Horndeski\} and \code{gravity_model_screening} = true & Combine the screened solution with the linear solution to get the correct linear evolution on large scales \\[2pt]
\param{gravity_model_screening_linear_scale_hmpc}\newline{\footnotesize double, \code{0.05}} & \code{gravity_model} in \{mbeta, f(R), Symmetron, Dilaton, DGP, CubicGalileon, Horndeski\} and \code{gravity_model_screening} = true & Use the linear solution for k < k\emph{ and the screened solution for k > k} (h/Mpc) \\[2pt]
\param{gravity_model_screening_efficiency}\newline{\footnotesize double, \code{1}} & \code{gravity_model} in \{mbeta, f(R), Symmetron, Dilaton, DGP, CubicGalileon, Horndeski\} and \code{gravity_model_screening} = true & Screening efficiency (1 is standard). Increases or decreases the amount of screening (it scales \code{Phi_crit}, or for DGP the density in the Vainshtein factor) \\[2pt]
\param{gravity_model_fofr_exact_solution}\newline{\footnotesize bool, \code{false}} & \code{gravity_model} = f(R) & Solve the exact f(R) equation with the multigrid solver (slow, mainly for testing) \\[2pt]
\param{gravity_model_fofr_multigrid_method}\newline{\footnotesize string, \code{"sqrt"}} & \code{gravity_model} = f(R) and \code{gravity_model_fofr_exact_solution} = true & The field the multigrid solver solves for. sqrt: u = sqrt(|\code{f_R}|) with the equation in a cell solved exactly (a cubic for n = 1), robust for any \code{f_R0}, n and redshift. log: ln|\code{f_R}| with Newton steps (the original method; can give NaN for small \code{f_R0} or at high redshift). Options: \code{sqrt}, \code{log} \\[2pt]
\param{gravity_model_symmetron_exact_solution}\newline{\footnotesize bool, \code{false}} & \code{gravity_model} = Symmetron & Solve the exact symmetron equation with the multigrid solver (slow, mainly for testing) \\[2pt]
\param{gravity_model_dilaton_exact_solution}\newline{\footnotesize bool, \code{false}} & \code{gravity_model} = Dilaton & Solve the exact dilaton equation with the multigrid solver (slow, mainly for testing) \\[2pt]
\param{gravity_model_dgp_exact_solution}\newline{\footnotesize bool, \code{false}} & \code{gravity_model} = DGP & Solve the exact DGP equation with the multigrid solver (slow, mainly for testing; can be hard to get to converge) \\[2pt]
\param{gravity_model_cubic_galileon_exact_solution}\newline{\footnotesize bool, \code{false}} & \code{gravity_model} = CubicGalileon & Solve the exact cubic Galileon equation (that of DGP with the Galileon's \code{r_c}(a) and beta(a)) with the multigrid solver (slow, mainly for testing) \\[2pt]
\param{gravity_model_horndeski_exact_solution}\newline{\footnotesize bool, \code{false}} & \code{gravity_model} = Horndeski & Solve the exact field equation (that of DGP with the table's \code{r_c}(a) and beta(a)) with the multigrid solver (slow, mainly for testing; needs \code{r_c}\textasciicircum{}2 > 0) \\[2pt]
\param{multigrid_nsweeps}\newline{\footnotesize int, \code{3}} & solving the exact equation of the gravity model & Multigrid solver: Newton-Gauss-Seidel sweeps on every level \\[2pt]
\param{multigrid_nsweeps_first_step}\newline{\footnotesize int, \code{20}} & solving the exact equation of the gravity model & Multigrid solver: sweeps on the first step (increase if the solver fails to converge) \\[2pt]
\param{multigrid_solver_residual_convergence}\newline{\footnotesize double, \code{1e-06}} & solving the exact equation of the gravity model & Multigrid solver: converged when the residual is below this \\[2pt]
\param{multigrid_convergence_criterion}\newline{\footnotesize string, \code{"residual"}} & solving the exact equation of the gravity model & Multigrid solver: converged when the rms residual is below \code{multigrid_solver_residual_convergence} (residual) or below it times the rms of the linear source of the equation (relative; the residual scale changes by orders of magnitude over a simulation). Options: \code{residual}, \code{relative} \\[2pt]
\param{multigrid_coarse_restriction}\newline{\footnotesize string, \code{"sharpk"}} & solving the exact equation of the gravity model & Multigrid solver with \code{multigrid_coarse_factor} > 1: the density on the coarse grid is its Fourier modes below the coarse Nyquist frequency (sharpk; delta < -1 set to -1) or a CIC deposit of the force-grid cells (cic; conserves the mass, delta >= -1). Options: \code{sharpk}, \code{cic} \\[2pt]
\param{multigrid_start_redshift}\newline{\footnotesize double, \code{-1}} & solving the exact equation of the gravity model & Multigrid solver: solve the field equation only below this redshift; above it the fifth force is set by \code{multigrid_force_before_start.} < 0: at all redshifts \\[2pt]
\param{multigrid_force_before_start}\newline{\footnotesize string, \code{"linear"}} & solving the exact equation of the gravity model & Multigrid solver with \code{multigrid_start_redshift} >= 0: the fifth force above that redshift is the linear one (linear; no screening) or none (GR; the growth factors, IC and COLA frame still have the modified gravity). Options: \code{linear}, \code{GR} \\[2pt]
\param{multigrid_coarse_factor}\newline{\footnotesize int, \code{1}} & solving the exact equation of the gravity model & Multigrid solver: solve the field equation on a grid coarser by this factor than the force grid (the density and the solution go between the grids as their Fourier modes below the coarse Nyquist frequency). 1: on the force grid \\[2pt]
\end{longtable}}

\section{Primordial power-spectrum}

{\small\begin{longtable}{@{}>{\raggedright\arraybackslash}p{4.3cm}>{\raggedright\arraybackslash}p{2.8cm}>{\raggedright\arraybackslash}p{8.1cm}@{}}
\toprule Parameter (type, default) & Used when & Description \\ \midrule
\endfirsthead
\toprule Parameter (type, default) & Used when & Description \\ \midrule
\endhead
\bottomrule
\endlastfoot
\param{cosmology_As}\newline{\footnotesize double, \code{2e-09}} & always & Primordial amplitude \\[2pt]
\param{cosmology_ns}\newline{\footnotesize double, \code{1}} & always & Spectral index \\[2pt]
\param{cosmology_kpivot_mpc}\newline{\footnotesize double, \code{0.05}} & always & Pivot scale (1/Mpc) \\[2pt]
\end{longtable}}

\section{Particles}

{\small\begin{longtable}{@{}>{\raggedright\arraybackslash}p{4.3cm}>{\raggedright\arraybackslash}p{2.8cm}>{\raggedright\arraybackslash}p{8.1cm}@{}}
\toprule Parameter (type, default) & Used when & Description \\ \midrule
\endfirsthead
\toprule Parameter (type, default) & Used when & Description \\ \midrule
\endhead
\bottomrule
\endlastfoot
\param{particle_Npart_1D}\newline{\footnotesize int, \emph{required}} & always & Number of CDM+baryon particles per dimension \\[2pt]
\param{particle_baryons}\newline{\footnotesize bool, \code{false}} & always & Evolve CDM and baryons as two sets of particles (\code{particle_Npart_1D}\textasciicircum{}3 of each, IC from the CDM and baryon transfer functions at zini). false: one set of CDM+baryon particles \\[2pt]
\param{particle_neutrinos}\newline{\footnotesize bool, \code{false}} & always & Massive neutrinos as particles with the delta-f method (Elbers et al. 2021): thermal Fermi-Dirac velocities, relativistic equations of motion, the particles only carry the perturbation. Replaces \code{force_linear_massive_neutrinos} (needs \code{cosmology_OmegaMNu} > 0, \code{ic_type_of_input} = camb, class, internal or tables, \code{timestep_method} = Quinn, no COLA) \\[2pt]
\param{particle_neutrino_Npart_1D}\newline{\footnotesize int, \code{0}} & \code{particle_neutrinos} = true & \code{particle_neutrinos}: number of neutrino particles per dimension (0: \code{particle_Npart_1D}) \\[2pt]
\param{particle_neutrino_deltaf}\newline{\footnotesize bool, \code{true}} & \code{particle_neutrinos} = true & \code{particle_neutrinos}: use the delta-f weights (false: plain particles, much noisier; for testing) \\[2pt]
\param{particle_baryon_lattice}\newline{\footnotesize string, \code{"same"}} & \code{particle_baryons} = true & \code{particle_baryons}: start the baryons from the same lattice as the CDM (same, recommended: the baryon-CDM difference then follows linear theory on all scales) or from a lattice offset by half a particle spacing (offset: gives a spurious baryon-CDM difference on small scales). Options: \code{same}, \code{offset} \\[2pt]
\param{particle_baryon_relative_mode}\newline{\footnotesize string, \code{"displacement"}} & \code{particle_baryons} = true & \code{particle_baryons}: how the baryon-CDM relative density is set up. displacement: each species is displaced by its own linear field (\code{T_i/T_cb}) and has its own velocity. mass: both species follow the total CDM+baryon flow (same displacements and velocities) and the relative density is carried by the particle masses \code{m_i} = 1 + \code{eps_i}(q) with \code{eps_i} = \code{delta_cb}(zini) (\code{T_i} - \code{T_cb})(z=0) / \code{T_cb}(zini), the linear relative mode today (Hahn, Rampf \& Uhlemann 2021; exact at z = 0 in linear theory, the relative mode is then constant in time). \code{mass_ini}: common displacements, masses with \code{eps_i} = \code{T_i} / \code{T_cb} - 1 at zini and each species keeps its own velocity (the relative mode then grows from zini). The mass modes need \code{particle_baryon_lattice} = same. Options: \code{displacement}, \code{mass}, \code{mass_ini} \\[2pt]
\param{particle_allocation_factor}\newline{\footnotesize double, \code{1.5}} & always & Allocate this factor more particles than the average per task (particles move between tasks) \\[2pt]
\param{particle_sort_every_nsteps}\newline{\footnotesize int, \code{0}} & always & Reorder the particles in memory by position (blocks of 4\textasciicircum{}3 force-grid cells) every this many steps, which makes the density assignment and the force interpolation faster when the particles have moved far from their initial (lattice) order. 0: never. Changes the results only by round-off (the order of the sums) and the order of the particles in the output \\[2pt]
\end{longtable}}

\section{Time-stepping}

{\small\begin{longtable}{@{}>{\raggedright\arraybackslash}p{4.3cm}>{\raggedright\arraybackslash}p{2.8cm}>{\raggedright\arraybackslash}p{8.1cm}@{}}
\toprule Parameter (type, default) & Used when & Description \\ \midrule
\endfirsthead
\toprule Parameter (type, default) & Used when & Description \\ \midrule
\endhead
\bottomrule
\endlastfoot
\param{timestep_nsteps}\newline{\footnotesize int list, \emph{required}} & always & Number of steps between the outputs (one per output redshift). If only one number it is the total number of steps \\[2pt]
\param{timestep_algorithm}\newline{\footnotesize string, \code{"KDK"}} & always & The time-stepping algorithm (only KDK is implemented). Options: \code{KDK} \\[2pt]
\param{timestep_method}\newline{\footnotesize string, \code{"Quinn"}} & always & The time-stepping method. Quinn: standard leapfrog. Tassev: modified leapfrog for COLA. FastPM: leapfrog that follows linear growth exactly (Feng et al. 2016). BullFrog: follows 2LPT growth after every step (Rampf, List \& Hahn 2024), not for COLA. Options: \code{Quinn}, \code{Tassev}, \code{FastPM}, \code{BullFrog} \\[2pt]
\param{timestep_cola_nLPT}\newline{\footnotesize double, \code{-2.5}} & always & The nLPT parameter of the Tassev time-stepping \\[2pt]
\param{timestep_scalefactor_spacing}\newline{\footnotesize string, \code{"linear"}} & always & Spacing of the steps in the scale factor (powerlaw: uniform in a\textasciicircum{}\code{timestep_spacing_power}, growthfactor: uniform in the linear growth factor D(a), as in Rampf, List \& Hahn 2024). Options: \code{linear}, \code{logarithmic}, \code{powerlaw}, \code{growthfactor} \\[2pt]
\param{timestep_spacing_power}\newline{\footnotesize double, \code{1}} & \code{timestep_scalefactor_spacing} = powerlaw & The steps are uniform in a\textasciicircum{}power \\[2pt]
\param{timestep_scaledependent_kick}\newline{\footnotesize bool, \code{false}} & \code{timestep_method} = FastPM & Experimental: multiply the force of each kick by a factor per k such that the linear growing mode of a scale-dependent growth D(k, a) (\code{G_eff}(k, a), massive neutrinos) is followed exactly for any step size (Zel'dovich-consistent at every k; the drift is unchanged and the velocities between the kicks are effective ones). Nothing changes without scale-dependent growth (see ScaleDependentKick.h) \\[2pt]
\end{longtable}}

\section{Initial conditions}

{\small\begin{longtable}{@{}>{\raggedright\arraybackslash}p{4.3cm}>{\raggedright\arraybackslash}p{2.8cm}>{\raggedright\arraybackslash}p{8.1cm}@{}}
\toprule Parameter (type, default) & Used when & Description \\ \midrule
\endfirsthead
\toprule Parameter (type, default) & Used when & Description \\ \midrule
\endhead
\bottomrule
\endlastfoot
\param{ic_initial_redshift}\newline{\footnotesize double, \emph{required}} & always & The initial redshift of the simulation \\[2pt]
\param{ic_random_seed}\newline{\footnotesize int, \emph{required}} & always & The random seed \\[2pt]
\param{ic_random_generator}\newline{\footnotesize string, \code{"MT19937"}} & always & The random number generator (GSL is \code{gsl_rng_ranlxd1} as used in 2LPTic, for comparisons). Options: \code{GSL}, \code{MT19937} \\[2pt]
\param{ic_random_field_backend}\newline{\footnotesize string, \code{"ngenic"}} & always & Fourier mode sampling: ngenic preserves legacy line seeds and IC; philox uses independent mode counters. Options: \code{ngenic}, \code{philox} \\[2pt]
\param{ic_nmesh}\newline{\footnotesize int, \emph{required}} & always & The grid size used to generate the IC \\[2pt]
\param{ic_lattice_offset}\newline{\footnotesize string, \code{"fiducial"}} & always & Offset of the initial particle lattice from the nodes of the force grid. fiducial: half a force grid cell if \code{force_nmesh} > \code{particle_Npart_1D}, otherwise none. With particles exactly on the nodes of a finer grid, CIC responds to the size of the displacement and gives large spurious forces (see the force grid section of EVOLUTION.md). \code{quarter_force_cell}: a quarter of a force cell (with \code{force_interlacing_order} = 2 no copy of the lattice is then on the nodes). Options: \code{fiducial}, \code{none}, \code{half_force_cell}, \code{quarter_force_cell} \\[2pt]
\param{ic_type_of_input}\newline{\footnotesize string, \emph{required}} & always & The input: powerspectrum (file with [k (h/Mpc), P(k) (Mpc/h)\textasciicircum{}3]), transferfunction (file with [k (h/Mpc), T(k) Mpc\textasciicircum{}2]), camb (a transfer info file listing T(k,z) files from CAMB; the old name transferinfofile is also accepted; the columns are read by name from the header, or by position: 13 CAMB, 9 AxionCAMB, 7 old CAMB / CLASS with format=camb), class (native CLASS output: \code{ic_input_filename} is the root of the run, with format = class, output = mTk,vTk, \code{write_background} = yes), internal (FML's Einstein-Boltzmann solver run with the cosmology of the simulation, \code{cosmology_model} included; the old name \code{internal_transfer} is also accepted) or tables (the tables of the internal solver in one file, made with ./nbody --write-transfers parameterfile.lua file; checked against the cosmology of the simulation). See FML/TransferProviders and \code{INITIAL_CONDITIONS.md.} Options: \code{powerspectrum}, \code{transferfunction}, \code{camb}, \code{class}, \code{internal}, \code{tables}, \code{internal_transfer}, \code{transferinfofile} \\[2pt]
\param{ic_input_gauge}\newline{\footnotesize string, \code{"synchronous"}} & \code{ic_type_of_input} = class & The gauge CLASS was run in (newtonian is converted to the CDM-synchronous gauge the IC need; it needs output = mTk,vTk). Options: \code{synchronous}, \code{newtonian} \\[2pt]
\param{ic_input_filename}\newline{\footnotesize string, \code{""}} & always & The input file (needed except for \code{ic_type_of_input} = internal, where it is an optional .ini file of the Einstein-Boltzmann solver with its precision, recombination and reionization settings: see \code{ic_internal_solver_parameters}) \\[2pt]
\param{ic_internal_solver_parameters}\newline{\footnotesize string list, \texttt{\{\}}} & \code{ic_type_of_input} = internal & Settings of the Einstein-Boltzmann solver as "name=value", e.g. \{"Yp=0.245", "\code{z_reion}=8", "precision=fast"\}, applied after those of the .ini file \code{ic_input_filename} (the names of Cosmology/EinsteinBoltzmannSolver/params/default.ini). The cosmology (h, OmegaB, ..., model and its parameters, \code{A_s}, \code{n_s}, \code{k_pivot}, the neutrino masses, preset), \code{alpha_s} and the k-grid (\code{pert_k_min}, \code{pert_k_max}, \code{pert_n_k}, \code{pert_dk_max}) are set by the simulation and are an error here and in the .ini file. See \code{INITIAL_CONDITIONS.md} \\[2pt]
\param{ic_axion_fraction}\newline{\footnotesize double, \code{0}} & \code{ic_type_of_input} = internal or tables & Ultralight axions in the IC only: the share of the dark matter (\code{cosmology_OmegaCDM}) that is an axion of mass 10\textasciicircum{}\code{ic_axion_log10_m} eV (the Einstein-Boltzmann solver's model axion, as axionCAMB). The particles are all the dark matter and start from the CDM + axion transfer functions (the axion's suppression of the small-scale power); the simulation itself is LCDM (\code{cosmology_model} = LCDM), so the axion's pressure after the start is neglected. 0 = off. Needs an axion that is a w = 0 fluid before the start (m >\textasciitilde{} 1e-26 eV for \code{z_ini} \textasciitilde{} 50). See \code{INITIAL_CONDITIONS.md} \\[2pt]
\param{ic_axion_log10_m}\newline{\footnotesize double, \code{-22}} & \code{ic_type_of_input} = internal or tables & log10 of the axion mass in eV (the solver's \code{ax_log10_m}; with \code{ic_axion_fraction} > 0) \\[2pt]
\param{ic_axion_switch_mH}\newline{\footnotesize double, \code{0}} & \code{ic_type_of_input} = internal or tables & The switch of the axion to its oscillation-averaged fluid (the solver's \code{ax_switch_mH}; with \code{ic_axion_fraction} > 0): 0 = once it oscillates fast (m/H \textasciitilde{} 125), 3 = at m = 3H as axionCAMB \\[2pt]
\param{ic_input_redshift}\newline{\footnotesize double, \emph{required}} & always & The redshift of the input P(k) / T(k) \\[2pt]
\param{ic_fix_amplitude}\newline{\footnotesize bool, \code{true}} & always & Fix the amplitude of the modes when generating the random field \\[2pt]
\param{ic_reverse_phases}\newline{\footnotesize bool, \code{false}} & always & Reverse the phases (for paired simulations) \\[2pt]
\param{ic_random_field_type}\newline{\footnotesize string, \code{"gaussian"}} & always & Type of IC: gaussian, nongaussian, \code{read_particles} (read a GADGET file and use that, reconstructing the LPT fields if COLA) or \code{read_phases} (take the phases from a GADGET file). Options: \code{gaussian}, \code{nongaussian}, \code{read_particles}, \code{read_phases} \\[2pt]
\param{ic_LPT_order}\newline{\footnotesize int, \code{2}} & always & The LPT order used to make the IC \\[2pt]
\param{ic_LPT_dealiasing}\newline{\footnotesize bool, \code{false}} & always & Dealias the nonlinear 2LPT and 3LPT sources: compute them on a 3N/2 (2LPT) or 2N (3LPT) mesh so that the products of the modes do not alias onto the grid (LPT/Dealiasing.h). Costs (3/2)\textasciicircum{}3 or 8 times the memory of an IC grid while the IC are made \\[2pt]
\param{ic_sigma8_normalization}\newline{\footnotesize bool, \code{false}} & always & Normalize P(k) to a given sigma8 (otherwise use the input) \\[2pt]
\param{ic_sigma8_redshift}\newline{\footnotesize double, \code{0}} & \code{ic_sigma8_normalization} = true & The redshift of the sigma8 value \\[2pt]
\param{ic_sigma8}\newline{\footnotesize double, \code{0.8}} & \code{ic_sigma8_normalization} = true & The sigma8 value to normalize to \\[2pt]
\param{ic_use_gravity_model_GR}\newline{\footnotesize bool, \code{false}} & always & For modified gravity: the input P(k) is for LCDM, use GR growth factors to scale it to the initial redshift (same IC as the LCDM simulation) \\[2pt]
\param{ic_fnl_type}\newline{\footnotesize string, \emph{required}} & \code{ic_random_field_type} = nongaussian & Type of non-gaussianity. Options: \code{local}, \code{equilateral}, \code{orthogonal} \\[2pt]
\param{ic_fnl}\newline{\footnotesize double, \emph{required}} & \code{ic_random_field_type} = nongaussian & The value of fNL \\[2pt]
\param{ic_fnl_redshift}\newline{\footnotesize double, \emph{required}} & \code{ic_random_field_type} = nongaussian & The redshift of the potential Phi = (3/5) g(z) zeta (g = D/a) that fNL multiplies: 0 is the usual large-scale-structure convention, a matter-era redshift (e.g. 30) the CMB convention \\[2pt]
\param{ic_reconstruct_gadgetfilepath}\newline{\footnotesize string, \emph{required}} & \code{ic_random_field_type} in \{\code{read_particles}, \code{read_phases}\} & Path to the GADGET files without the file number: binary (path.0, path.1, ...) or HDF5 (path.hdf5 or path.0.hdf5, ...), detected automatically \\[2pt]
\param{ic_reconstruct_assigment_method}\newline{\footnotesize string, \code{"CIC"}} & \code{ic_random_field_type} in \{\code{read_particles}, \code{read_phases}\} & Density assignment method used to reconstruct the density field. Options: \code{NGP}, \code{CIC}, \code{TSC}, \code{PCS}, \code{PQS} \\[2pt]
\param{ic_reconstruct_smoothing_filter}\newline{\footnotesize string, \code{"sharpk"}} & \code{ic_random_field_type} in \{\code{read_particles}, \code{read_phases}\} & Filter used to remove small-scale modes from the reconstructed density field. Options: \code{tophat}, \code{sharpk}, \code{gaussian} \\[2pt]
\param{ic_reconstruct_dimless_smoothing_scale}\newline{\footnotesize double, \code{0}} & \code{ic_random_field_type} in \{\code{read_particles}, \code{read_phases}\} & Smoothing scale R/boxsize of the reconstructed density field \\[2pt]
\param{ic_reconstruct_interlacing}\newline{\footnotesize bool, \code{false}} & \code{ic_random_field_type} in \{\code{read_particles}, \code{read_phases}\} & Use interlacing when reconstructing the density field \\[2pt]
\param{ic_reconstruct_exact}\newline{\footnotesize bool, \code{false}} & \code{ic_random_field_type} in \{\code{read_particles}, \code{read_phases}\} & Reconstruct the LPT fields exactly from the positions and velocities (always works for 1LPT IC; for 2LPT IC the particles must come from a regular grid) \\[2pt]
\end{longtable}}

\section{Forces}

{\small\begin{longtable}{@{}>{\raggedright\arraybackslash}p{4.3cm}>{\raggedright\arraybackslash}p{2.8cm}>{\raggedright\arraybackslash}p{8.1cm}@{}}
\toprule Parameter (type, default) & Used when & Description \\ \midrule
\endfirsthead
\toprule Parameter (type, default) & Used when & Description \\ \midrule
\endhead
\bottomrule
\endlastfoot
\param{force_nmesh}\newline{\footnotesize int, \emph{required}} & always & Grid size used to compute the PM forces \\[2pt]
\param{force_density_assignment_method}\newline{\footnotesize string, \code{"CIC"}} & always & Density assignment method (not NGP: its force only changes when a particle crosses a cell boundary, so near the initial lattice it is zero or jumps by order one, see EVOLUTION.md). Options: \code{CIC}, \code{TSC}, \code{PCS}, \code{PQS} \\[2pt]
\param{force_greens_function_kernel}\newline{\footnotesize string, \code{"fiducial"}} & always & Kernel for D\textasciicircum{}2 when solving the Poisson equation (fiducial = continuous). Options: \code{fiducial}, \code{continuous}, \code{discrete_2pt}, \code{discrete_4pt} \\[2pt]
\param{force_gradient_kernel}\newline{\footnotesize string, \code{"fiducial"}} & always & Kernel for D when computing the forces. fiducial = continuous if \code{force_nmesh} = \code{particle_Npart_1D}, \code{discrete_4pt} if \code{force_nmesh} is up to 3x larger and \code{discrete_2pt} if it is even larger (the continuous kernel gives spurious forces from the particle lattice when the force grid is finer than it). Options: \code{fiducial}, \code{continuous}, \code{discrete_2pt}, \code{discrete_4pt} \\[2pt]
\param{force_density_method}\newline{\footnotesize string, \code{"assignment"}} & always & How the density field for the force is computed. assignment: particles to grid with \code{force_density_assignment_method.} nufft: the exact Fourier modes of the particle distribution (non-uniform FFT; no window function or aliasing). sheet: spectral sheet interpolation, R\textasciicircum{}NDIM tracers per particle placed by Fourier interpolating the displacement field x - q, assigned with \code{force_density_assignment_method.} \code{sheet_nufft}: sheet interpolation with the exact Fourier modes of the tracers (NUFFT). The force is always interpolated back to the particles with \code{force_density_assignment_method.} Options: \code{assignment}, \code{nufft}, \code{sheet}, \code{sheet_nufft} \\[2pt]
\param{force_interlacing_order}\newline{\footnotesize int, \code{1}} & always & Force interlacing: compute the force for this many copies of the particles shifted by a fraction of a cell relative to the grid and average (at the particles). 1: off, 2: shift (1/2,1/2,1/2), 4: the half-cell face diagonals, 8: all half-cell corners, 16, 32: the half-cell corners refined along the diagonal, n\textasciicircum{}3 (27, 64, ...): a sub-lattice of n points per cell side, any other order below 64 (3, 5, 6, 7, 9, ...): a rank-1 lattice j g / M (7 and 9 cancel about as much as 16, 13 more). Cancels the leading aliasing of the density assignment and the force interpolation (the images at k + 2 \code{k_Ny} n with only one nonzero \code{n_i}, \code{n_i} = 1) at the cost of this many force calculations. 2 and 4 also cancel some of the others but leave n = (2,0,0) (so does 8; the rank-1 lattices cancel it); 16 and 64 leave (4,0,0), 27 leaves (3,0,0) (see \code{generate_interlacing_shifts}) \\[2pt]
\param{force_deconvolve_window}\newline{\footnotesize bool, \code{false}} & always & Divide the force by W(k)\textasciicircum{}2 where W is the window function of \code{force_density_assignment_method} (corrects for the smoothing of the assignment and the interpolation; by W(k) if one of them is a NUFFT: \code{force_density_method} = nufft or \code{sheet_nufft}, or \code{force_interpolation_method} = nufft; not at all if both are). Unstable unless the aliasing is suppressed (\code{force_interlacing_order} > 1, see EVOLUTION.md) \\[2pt]
\param{force_lattice_greens_function}\newline{\footnotesize bool, \code{false}} & always & Multiply the Green's function by the lattice-calibrated factor (1 - E)/(T - E), with T(k) the linear response of the PM force for particles on the initial lattice, computed in closed form from the assignment, kernels, grid factor, lattice offset, interlacing and deconvolution. The linear force is then exact up to the particle Nyquist frequency while the particles are close to the lattice. Needs \code{force_nmesh} = an integer times \code{particle_Npart_1D}, assignment density and interpolation, IC on a lattice and no CIC copy of the lattice on the grid nodes (use \code{ic_lattice_offset}). See EVOLUTION.md \\[2pt]
\param{force_lattice_greens_until_redshift}\newline{\footnotesize double, \code{-1}} & \code{force_lattice_greens_function} = true & Use the lattice-calibrated Green's function for steps with z > this (-1: always) \\[2pt]
\param{force_lattice_greens_taper_kmin}\newline{\footnotesize double, \code{-1}} & \code{force_lattice_greens_function} = true & The factor is applied fully below this wavenumber of the particle mode (units of the particle Nyquist frequency) and tapered to 1 at the particle Nyquist frequency. -1: fiducial (1, no taper, if \code{force_nmesh} = \code{particle_Npart_1D}, otherwise 0.7: on a finer grid the untapered factor is unstable over many steps) \\[2pt]
\param{force_lattice_greens_cap}\newline{\footnotesize double, \code{-1}} & \code{force_lattice_greens_function} = true & Clip the factor to [1/cap, cap]. -1: fiducial (no clipping if \code{force_nmesh} = \code{particle_Npart_1D}, otherwise 3) \\[2pt]
\param{force_early_method}\newline{\footnotesize string, \code{"none"}} & always & A more accurate force at early times (when the particles are still close to the lattice and CIC is inaccurate): nufft (NUFFT density and interpolation) or interlaced (\code{force_early_interlacing_order} shifted grids with \code{force_early_density_assignment_method} and window deconvolution) for steps at z > \code{force_early_until_redshift}, then the normal force settings. none: off. Options: \code{none}, \code{nufft}, \code{interlaced} \\[2pt]
\param{force_early_until_redshift}\newline{\footnotesize double, \code{10}} & \code{force_early_method} in \{nufft, interlaced\} & Use the early force for z > this \\[2pt]
\param{force_early_interlacing_order}\newline{\footnotesize int, \code{2}} & \code{force_early_method} = interlaced & Number of shifted grids at early times \\[2pt]
\param{force_early_density_assignment_method}\newline{\footnotesize string, \code{"TSC"}} & \code{force_early_method} = interlaced & Assignment kernel at early times (not NGP). Options: \code{CIC}, \code{TSC}, \code{PCS}, \code{PQS} \\[2pt]
\param{force_interpolation_method}\newline{\footnotesize string, \code{"fiducial"}} & always & How the force is interpolated to the particles. assignment: with the \code{force_density_assignment_method} kernel. nufft: exact evaluation of the Fourier series of the force at the particles (non-uniform FFT). fiducial: nufft if \code{force_density_method} = nufft (then assignment and interpolation are both exact, so there is no self-force), otherwise assignment. Options: \code{fiducial}, \code{assignment}, \code{nufft} \\[2pt]
\param{force_nufft_epsilon}\newline{\footnotesize double, \code{1e-05}} & \code{force_density_method} in \{nufft, \code{sheet_nufft}\} or \code{force_interpolation_method} = nufft or \code{force_early_method} = nufft & Relative accuracy of the NUFFT (\code{force_density_method} = nufft or \code{sheet_nufft}, \code{force_interpolation_method} = nufft or \code{force_early_method} = nufft) \\[2pt]
\param{force_sheet_refinement}\newline{\footnotesize int, \code{2}} & \code{force_density_method} in \{sheet, \code{sheet_nufft}\} & R: the number of tracers per particle per dimension for sheet interpolation \\[2pt]
\param{force_linear_massive_neutrinos}\newline{\footnotesize bool, \code{false}} & always & Include massive neutrinos in the density field using linear theory (requires \code{ic_type_of_input} = camb, class, internal or tables) \\[2pt]
\param{force_neutrino_energy_density}\newline{\footnotesize bool, \code{false}} & always & Source the Poisson equation with the energy density of the massive neutrinos instead of their mass density: \code{f_nu} <E/m>(a) \code{delta_nu} with <E/m> = <sqrt(1 + (q/am)\textasciicircum{}2)> over the Fermi-Dirac distribution (both for \code{force_linear_massive_neutrinos}, where \code{delta_nu} from CAMB/CLASS is already the energy-density contrast, and for \code{particle_neutrinos}, where each particle is weighted by its E/m and the neutrino P(k) is then the energy-density contrast). Matters at high redshift for light neutrinos (<E/m> = 1.33 at z = 50 for 0.033 eV) \\[2pt]
\param{force_nbody_gauge}\newline{\footnotesize bool, \code{false}} & always & Run in the N-body gauge: the initial conditions in the N-body gauge and the linear source from radiation, neutrinos and GR added to the density field. With \code{ic_type_of_input} = internal both come from the Einstein-Boltzmann solver; with camb from the tables made with CLASS by \code{scripts/make_nbody_gauge_transfer.py} \\[2pt]
\param{force_nbody_gauge_filename}\newline{\footnotesize string, \emph{required}} & \code{force_nbody_gauge} = true and \code{ic_type_of_input} = camb & The table \code{T_extra}(k,z) (\code{nbgauge_extra.txt} from \code{scripts/make_nbody_gauge_transfer.py}) \\[2pt]
\param{force_use_finite_difference_force}\newline{\footnotesize bool, \code{false}} & always & Experimental: compute forces with finite differences of the potential \\[2pt]
\param{force_finite_difference_stencil_order}\newline{\footnotesize int, \code{4}} & always & Experimental: order of the finite difference stencil \\[2pt]
\end{longtable}}

\section{Lightcone}

{\small\begin{longtable}{@{}>{\raggedright\arraybackslash}p{4.3cm}>{\raggedright\arraybackslash}p{2.8cm}>{\raggedright\arraybackslash}p{8.1cm}@{}}
\toprule Parameter (type, default) & Used when & Description \\ \midrule
\endfirsthead
\toprule Parameter (type, default) & Used when & Description \\ \midrule
\endhead
\bottomrule
\endlastfoot
\param{lightcone}\newline{\footnotesize bool, \code{false}} & always & Make a lightcone \\[2pt]
\param{plc_pos_observer}\newline{\footnotesize double list, \texttt{\{0, 0, 0\}}} & lightcone = true & Position of the observer in units of the boxsize \\[2pt]
\param{plc_z_init}\newline{\footnotesize double, \code{0}} & lightcone = true & The redshift we start recording the lightcone \\[2pt]
\param{plc_z_finish}\newline{\footnotesize double, \code{0}} & lightcone = true & The redshift we stop recording the lightcone \\[2pt]
\param{plc_boundary_mpch}\newline{\footnotesize double, \code{20}} & lightcone = true & Boundary region around the shell to make sure we get all particles (Mpc/h) \\[2pt]
\param{plc_use_replicas}\newline{\footnotesize bool, \code{false}} & lightcone = true & Replicate the box to get the sky coverage we want \\[2pt]
\param{plc_ndim_rep}\newline{\footnotesize int, \code{0}} & lightcone = true & Number of dimensions with replicas in both directions (3D: 0 = octant, 3 = full sky) \\[2pt]
\param{plc_output_gadgetfile}\newline{\footnotesize bool, \code{false}} & lightcone = true & Output the lightcone particles in GADGET format \\[2pt]
\param{plc_output_asciifile}\newline{\footnotesize bool, \code{false}} & lightcone = true & Output the lightcone particles in ascii \\[2pt]
\param{plc_output_hdf5file}\newline{\footnotesize bool, \code{false}} & lightcone = true & Output the lightcone particles as GADGET HDF5 files (needs \code{USE_HDF5}; positions in Mpc/h, peculiar velocities in km/s) \\[2pt]
\param{plc_output_in_batches}\newline{\footnotesize bool, \code{false}} & lightcone = true & Output in batches to save memory \\[2pt]
\param{plc_relativistic_effects}\newline{\footnotesize bool, \code{false}} & lightcone = true & Store the potential Phi/c\textasciicircum{}2 at the lightcone particles and write it with the radial velocity, the cosmological redshift and the observed redshift 1 + \code{z_obs} = (1 + z)(1 + \code{v_r/c} - Phi/c\textasciicircum{}2) (Doppler + gravitational redshift; observer at rest with Phi = 0) to the HDF5 and ascii lightcone files. With \code{plc_make_onion_density_maps} also Healpix maps of Phi/c\textasciicircum{}2 at the shells. Phi from the Poisson equation with the density used for the force (the GR potential also for modified gravity models). Costs one extra FFT per step while recording the lightcone \\[2pt]
\param{plc_halos}\newline{\footnotesize bool, \code{false}} & lightcone = true & Friends-of-friends halos on the lightcone: at every step that records the lightcone, FoF on the particles near the lightcone shell of the step (with their true velocities); each halo is recorded once, when its center of mass crosses the lightcone (moving with the mean velocity of its particles, as the particles do), in \code{lightcone_<name>/halos_lightcone.txt} \\[2pt]
\param{plc_halos_nmin_per_halo}\newline{\footnotesize int, \code{20}} & lightcone = true and \code{plc_halos} = true & Minimum number of particles per lightcone halo \\[2pt]
\param{plc_halos_linking_length}\newline{\footnotesize double, \code{0.2}} & lightcone = true and \code{plc_halos} = true & Linking length of the lightcone halos in units of the mean particle separation \\[2pt]
\param{plc_halos_buffer_mpch}\newline{\footnotesize double, \code{0}} & lightcone = true and \code{plc_halos} = true & FoF buffer region from neighbor tasks for the lightcone halos (Mpc/h; larger than the largest halo). 0: the larger of 3 Mpc/h and 10 linking lengths \\[2pt]
\param{plc_halos_shell_only}\newline{\footnotesize bool, \code{true}} & lightcone = true and \code{plc_halos} = true & Link only the particles within \code{plc_boundary_mpch} (plus the largest displacement of the step) of the lightcone shell of the step (in any replica). false: all the particles every step (for tests; the same halos) \\[2pt]
\param{plc_make_onion_density_maps}\newline{\footnotesize bool, \code{false}} & lightcone = true & Make (Healpix) density maps delta(z, theta) \\[2pt]
\param{plc_nside}\newline{\footnotesize int, \code{256}} & lightcone = true and \code{plc_make_onion_density_maps} = true & Healpix nside of the maps (npix = 12 nside\textasciicircum{}2; at most 8192) \\[2pt]
\param{plc_da_maps}\newline{\footnotesize double, \code{0.02}} & lightcone = true and \code{plc_make_onion_density_maps} = true & Roughly the size of the bins in a for the maps \\[2pt]
\param{plc_use_chunkpix}\newline{\footnotesize bool, \code{false}} & lightcone = true and \code{plc_make_onion_density_maps} = true & Store the maps with chunkpix (useful for sparse maps) \\[2pt]
\param{plc_nside_chunks}\newline{\footnotesize int, \code{8}} & lightcone = true and \code{plc_make_onion_density_maps} = true and \code{plc_use_chunkpix} = true & Healpix nside of the chunks (at most 8192) \\[2pt]
\end{longtable}}

\section{Output}

{\small\begin{longtable}{@{}>{\raggedright\arraybackslash}p{4.3cm}>{\raggedright\arraybackslash}p{2.8cm}>{\raggedright\arraybackslash}p{8.1cm}@{}}
\toprule Parameter (type, default) & Used when & Description \\ \midrule
\endfirsthead
\toprule Parameter (type, default) & Used when & Description \\ \midrule
\endhead
\bottomrule
\endlastfoot
\param{output_folder}\newline{\footnotesize string, \emph{required}} & always & Folder to put the output in ("" is the current folder) \\[2pt]
\param{output_redshifts}\newline{\footnotesize double list, \emph{required}} & always & The output redshifts (decreasing) \\[2pt]
\param{output_particles}\newline{\footnotesize bool, \code{true}} & always & Write the particles at the output redshifts \\[2pt]
\param{output_fileformat}\newline{\footnotesize string, \code{"GADGET"}} & always & File format for the particles. GADGET: binary GADGET files. HDF5: GADGET HDF5 files (compile with \code{USE_HDF5}=true). Options: \code{GADGET}, \code{HDF5} \\[2pt]
\end{longtable}}

\section{Baryonification (snapshot post-processing)}

{\small\begin{longtable}{@{}>{\raggedright\arraybackslash}p{4.3cm}>{\raggedright\arraybackslash}p{2.8cm}>{\raggedright\arraybackslash}p{8.1cm}@{}}
\toprule Parameter (type, default) & Used when & Description \\ \midrule
\endfirsthead
\toprule Parameter (type, default) & Used when & Description \\ \midrule
\endhead
\bottomrule
\endlastfoot
\param{baryonification}\newline{\footnotesize bool, \code{false}} & always & Correct a temporary snapshot copy; 3D GR, no massive neutrinos, single CDM+baryon species; needs pofk=true \\[2pt]
\param{baryonification_model}\newline{\footnotesize string, \code{"schneider2019"}} & baryonification = true & Finite-support Schneider2019 profile family; illustrative, uncalibrated defaults \\[2pt]
\param{baryonification_output_particles}\newline{\footnotesize bool, \code{false}} & baryonification = true & Write a separate corrected snapshot; velocities are unchanged \\[2pt]
\param{baryonification_min_halo_mass_msunh}\newline{\footnotesize double, \code{1e+12}} & baryonification = true & Minimum host M200c (Msun/h) \\[2pt]
\param{baryonification_gas_mass_scale_msunh}\newline{\footnotesize double, \code{1e+14}} & baryonification = true & Gas-slope mass scale \code{M_c} (Msun/h) \\[2pt]
\param{baryonification_stellar_mass_scale_msunh}\newline{\footnotesize double, \code{2.5e+11}} & baryonification = true & Stellar-fraction pivot M1 (Msun/h) \\[2pt]
\param{baryonification_gas_mass_exponent}\newline{\footnotesize double, \code{0.3}} & baryonification = true & Exponent mu in beta=3-(\code{M_c/M200c})\textasciicircum{}mu \\[2pt]
\param{baryonification_gas_core_radius_over_r200}\newline{\footnotesize double, \code{0.1}} & baryonification = true & Gas core radius / R200c \\[2pt]
\param{baryonification_gas_ejection_radius_over_r200}\newline{\footnotesize double, \code{4}} & baryonification = true & Gas ejection radius / R200c \\[2pt]
\param{baryonification_truncation_radius_over_r200}\newline{\footnotesize double, \code{4}} & baryonification = true & NFW truncation radius / R200c \\[2pt]
\param{baryonification_stellar_radius_over_r200}\newline{\footnotesize double, \code{0.015}} & baryonification = true & Central stellar scale \code{R_h} / R200c \\[2pt]
\param{baryonification_stellar_amplitude}\newline{\footnotesize double, \code{0.09}} & baryonification = true & Stellar mass fraction at M1 \\[2pt]
\param{baryonification_stellar_total_exponent}\newline{\footnotesize double, \code{0.32}} & baryonification = true & Total stellar fraction power-law exponent \\[2pt]
\param{baryonification_stellar_central_exponent}\newline{\footnotesize double, \code{0.6}} & baryonification = true & Central stellar fraction power-law exponent \\[2pt]
\param{baryonification_relax_collisionless}\newline{\footnotesize bool, \code{true}} & baryonification = true & Apply collisionless contraction/expansion \\[2pt]
\param{baryonification_relaxation_amplitude}\newline{\footnotesize double, \code{0.3}} & baryonification = true & Dimensionless collisionless relaxation amplitude \\[2pt]
\param{baryonification_relaxation_exponent}\newline{\footnotesize double, \code{2}} & baryonification = true & Dimensionless collisionless relaxation exponent \\[2pt]
\param{baryonification_concentration_norm}\newline{\footnotesize double, \code{5}} & baryonification = true & Illustrative c200c normalization; not a fitted relation \\[2pt]
\param{baryonification_concentration_pivot_mass_msunh}\newline{\footnotesize double, \code{1e+14}} & baryonification = true & Concentration pivot mass (Msun/h) \\[2pt]
\param{baryonification_concentration_mass_slope}\newline{\footnotesize double, \code{-0.1}} & baryonification = true & c=c0\emph{(M200c/Mpivot)\textasciicircum{}slope}(1+z)\textasciicircum{}\code{redshift_power} \\[2pt]
\param{baryonification_concentration_redshift_power}\newline{\footnotesize double, \code{-1}} & baryonification = true & Power of (1+z) in the concentration relation \\[2pt]
\param{baryonification_fof_linking_length}\newline{\footnotesize double, \code{0.2}} & baryonification = true & FoF seed linking length / mean particle separation \\[2pt]
\param{baryonification_fof_min_particles}\newline{\footnotesize int, \code{20}} & baryonification = true & Minimum particles in a FoF seed halo \\[2pt]
\param{baryonification_fof_nmesh_max}\newline{\footnotesize int, \code{0}} & baryonification = true & Maximum FoF search mesh (0: automatic) \\[2pt]
\param{baryonification_fof_buffer_mpch}\newline{\footnotesize double, \code{3}} & baryonification = true & Halo/SO search buffer (comoving Mpc/h); larger than halo radii \\[2pt]
\param{baryonification_profile_radial_bins}\newline{\footnotesize int, \code{2048}} & baryonification = true & Number of positive profile-table radii \\[2pt]
\param{baryonification_min_radius_over_r200}\newline{\footnotesize double, \code{1e-06}} & baryonification = true & First positive profile radius / R200c \\[2pt]
\param{baryonification_max_radius_over_r200}\newline{\footnotesize double, \code{30}} & baryonification = true & Finite profile support / R200c; must be smaller than half the box \\[2pt]
\end{longtable}}

\section{Halo finding}

{\small\begin{longtable}{@{}>{\raggedright\arraybackslash}p{4.3cm}>{\raggedright\arraybackslash}p{2.8cm}>{\raggedright\arraybackslash}p{8.1cm}@{}}
\toprule Parameter (type, default) & Used when & Description \\ \midrule
\endfirsthead
\toprule Parameter (type, default) & Used when & Description \\ \midrule
\endhead
\bottomrule
\endlastfoot
\param{fof}\newline{\footnotesize bool, \code{false}} & always & Locate friends-of-friends halos at every output \\[2pt]
\param{fof_nmin_per_halo}\newline{\footnotesize int, \code{20}} & fof = true & Minimum number of particles per halo (of one species with \code{particle_baryons}) \\[2pt]
\param{fof_linking_length}\newline{\footnotesize double, \code{0.2}} & fof = true & The linking length in units of the mean particle separation (of one species with \code{particle_baryons}) \\[2pt]
\param{fof_nmesh_max}\newline{\footnotesize int, \code{0}} & fof = true & Maximum grid size used to bin particles to speed up the linking (0: the code chooses) \\[2pt]
\param{fof_buffer_length_mpch}\newline{\footnotesize double, \code{3}} & fof = true & Buffer region (larger than the largest halo) from neighbor tasks (Mpc/h). The code warns if halos reach its edge (they might then be cut off and found twice) \\[2pt]
\param{fof_output_fileformat}\newline{\footnotesize string, \code{"ascii"}} & fof = true & fof: the halo catalogue \code{halos_z<z>.txt} (ascii, one row per halo) or \code{halos_z<z>.hdf5} (HDF5, compile with \code{USE_HDF5}=true). Options: \code{ascii}, \code{HDF5} \\[2pt]
\param{fof_massfunction_nbins}\newline{\footnotesize int, \code{30}} & fof = true & fof: number of logarithmic bins of the mass function \\[2pt]
\param{fof_massfunction_Mmin}\newline{\footnotesize double, \code{1e+10}} & fof = true & fof: smallest mass of the mass function (Msun/h) \\[2pt]
\param{fof_massfunction_Mmax}\newline{\footnotesize double, \code{1e+16}} & fof = true & fof: largest mass of the mass function (Msun/h) \\[2pt]
\param{fof_so_definitions}\newline{\footnotesize string list, \texttt{\{\}}} & fof = true & fof: spherical overdensity masses around the FoF halos, e.g. \{"200m", "500c", "vir"\} (<Delta>m: Delta times the mean CDM+baryon density, <Delta>c: the critical density, vir: Bryan \& Norman 1998). Empty: none. The spheres must fit in \code{fof_buffer_length_mpch} \\[2pt]
\param{fof_relaxed}\newline{\footnotesize bool, \code{false}} & fof = true & fof: relaxed FoF instead of standard FoF (Dai et al. 2019; FastPM's RFOF; see FriendsOfFriends/RelaxedFoF.md): a linking length that is longer for small halos, and rejection of fake halos by their velocity dispersion. \code{fof_linking_length} and \code{fof_nmin_per_halo} are then not used \\[2pt]
\param{fof_relaxed_mode}\newline{\footnotesize string, \code{"fastpm"}} & fof = true and \code{fof_relaxed} = true & fastpm: the halo bins only, as FastPM (its parameters were fitted for this); paper: then shorter linking lengths (x 0.9 per pass) for the particles of rejected halos, as Algorithm 1 of the paper. Options: \code{fastpm}, \code{paper} \\[2pt]
\param{fof_relaxed_l1}\newline{\footnotesize double, \code{0.25}} & fof = true and \code{fof_relaxed} = true & Linking length of the smallest halos: l1 - a1 / (1 + z) (units of the mean particle separation; FastPM's defaults are fitted to TNG300-2-Dark for FastPM with a mean separation of 0.3 Mpc/h and 40 steps) \\[2pt]
\param{fof_relaxed_a1}\newline{\footnotesize double, \code{0.012}} & fof = true and \code{fof_relaxed} = true & See \code{fof_relaxed_l1} \\[2pt]
\param{fof_relaxed_l6}\newline{\footnotesize double, \code{0.24}} & fof = true and \code{fof_relaxed} = true & Linking length of the largest halos: max(l6 - a2 / (1 + z), \code{fof_relaxed_linking_length_min}); linear in the bin in between \\[2pt]
\param{fof_relaxed_a2}\newline{\footnotesize double, \code{0.06}} & fof = true and \code{fof_relaxed} = true & See \code{fof_relaxed_l6} \\[2pt]
\param{fof_relaxed_linking_length_min}\newline{\footnotesize double, \code{0.2}} & fof = true and \code{fof_relaxed} = true & Smallest linking length of the largest halos (units of the mean particle separation) \\[2pt]
\param{fof_relaxed_bin_edges}\newline{\footnotesize int list, \texttt{\{20, 40, 80, 160, 320\}}} & fof = true and \code{fof_relaxed} = true & Upper edges (particles) of the halo bins; the last bin has no upper edge \\[2pt]
\param{fof_relaxed_nmin}\newline{\footnotesize int, \code{8}} & fof = true and \code{fof_relaxed} = true & Minimum number of particles per halo (of one species with \code{particle_baryons}) \\[2pt]
\param{fof_relaxed_velocity_rejection}\newline{\footnotesize bool, \code{true}} & fof = true and \code{fof_relaxed} = true & Reject halos with a 3D velocity dispersion above r0(z) V0 (E(z) M / 1e15 Msun/h)\textasciicircum{}(1/3), V0 = 1100 km/s, r0 = b1 - b2 ln(1 + z) \\[2pt]
\param{fof_relaxed_b1}\newline{\footnotesize double, \code{7.02}} & fof = true and \code{fof_relaxed} = true & See \code{fof_relaxed_velocity_rejection} \\[2pt]
\param{fof_relaxed_b2}\newline{\footnotesize double, \code{6.025}} & fof = true and \code{fof_relaxed} = true & See \code{fof_relaxed_velocity_rejection.} With the defaults r0 is 0.40 at z = 2 and <= 0 for z > 2.2 (all halos rejected): they were fitted at z <= 2 \\[2pt]
\param{fof_relaxed_r0_min}\newline{\footnotesize double, \code{0}} & fof = true and \code{fof_relaxed} = true & Floor of r0(z) (0: none, as FastPM) \\[2pt]
\end{longtable}}

\section{Power-spectrum}

{\small\begin{longtable}{@{}>{\raggedright\arraybackslash}p{4.3cm}>{\raggedright\arraybackslash}p{2.8cm}>{\raggedright\arraybackslash}p{8.1cm}@{}}
\toprule Parameter (type, default) & Used when & Description \\ \midrule
\endfirsthead
\toprule Parameter (type, default) & Used when & Description \\ \midrule
\endhead
\bottomrule
\endlastfoot
\param{pofk}\newline{\footnotesize bool, \code{false}} & always & Compute the power-spectrum at every output \\[2pt]
\param{pofk_every_step}\newline{\footnotesize bool, \code{true}} & always & Write P(k) of the density field used for the force at every step (cheap; same grid and assignment as the force, the window is deconvolved): \code{pofk_<name>_cb_z<z>.txt} in \code{output_folder} \\[2pt]
\param{pofk_subtract_shotnoise}\newline{\footnotesize bool, \code{false}} & pofk = true or \code{pofk_every_step} = true & Subtract shot-noise (pofk and \code{pofk_every_step}) \\[2pt]
\param{pofk_nmesh}\newline{\footnotesize int, \emph{required}} & pofk = true & Grid size \\[2pt]
\param{pofk_interlacing}\newline{\footnotesize bool, \code{false}} & pofk = true & Use interlaced grids to reduce aliasing \\[2pt]
\param{pofk_density_assignment_method}\newline{\footnotesize string, \code{"CIC"}} & pofk = true & Density assignment method. Options: \code{NGP}, \code{CIC}, \code{TSC}, \code{PCS}, \code{PQS} \\[2pt]
\end{longtable}}

\section{Power-spectrum multipoles}

{\small\begin{longtable}{@{}>{\raggedright\arraybackslash}p{4.3cm}>{\raggedright\arraybackslash}p{2.8cm}>{\raggedright\arraybackslash}p{8.1cm}@{}}
\toprule Parameter (type, default) & Used when & Description \\ \midrule
\endfirsthead
\toprule Parameter (type, default) & Used when & Description \\ \midrule
\endhead
\bottomrule
\endlastfoot
\param{pofk_multipole}\newline{\footnotesize bool, \code{false}} & always & Compute the redshift space multipoles \code{P_ell}(k) at every output \\[2pt]
\param{pofk_multipole_nmesh}\newline{\footnotesize int, \emph{required}} & \code{pofk_multipole} = true & Grid size \\[2pt]
\param{pofk_multipole_interlacing}\newline{\footnotesize bool, \code{false}} & \code{pofk_multipole} = true & Use interlaced grids to reduce aliasing \\[2pt]
\param{pofk_multipole_subtract_shotnoise}\newline{\footnotesize bool, \code{false}} & \code{pofk_multipole} = true & Subtract shot-noise from P0 \\[2pt]
\param{pofk_multipole_ellmax}\newline{\footnotesize int, \code{4}} & \code{pofk_multipole} = true & Compute \code{P_ell} for ell = 0, 2, ..., ellmax \\[2pt]
\param{pofk_multipole_density_assignment_method}\newline{\footnotesize string, \code{"CIC"}} & \code{pofk_multipole} = true & Density assignment method. Options: \code{NGP}, \code{CIC}, \code{TSC}, \code{PCS}, \code{PQS} \\[2pt]
\end{longtable}}

\section{Bispectrum}

{\small\begin{longtable}{@{}>{\raggedright\arraybackslash}p{4.3cm}>{\raggedright\arraybackslash}p{2.8cm}>{\raggedright\arraybackslash}p{8.1cm}@{}}
\toprule Parameter (type, default) & Used when & Description \\ \midrule
\endfirsthead
\toprule Parameter (type, default) & Used when & Description \\ \midrule
\endhead
\bottomrule
\endlastfoot
\param{bispectrum}\newline{\footnotesize bool, \code{false}} & always & Compute the bispectrum at every output \\[2pt]
\param{bispectrum_nmesh}\newline{\footnotesize int, \emph{required}} & bispectrum = true & Grid size \\[2pt]
\param{bispectrum_nbins}\newline{\footnotesize int, \emph{required}} & bispectrum = true & Number of k-bins (memory and time scale as nbins grids and nbins\textasciicircum{}3 integrals) \\[2pt]
\param{bispectrum_density_assignment_method}\newline{\footnotesize string, \code{"CIC"}} & bispectrum = true & Density assignment method. Options: \code{NGP}, \code{CIC}, \code{TSC}, \code{PCS}, \code{PQS} \\[2pt]
\param{bispectrum_interlacing}\newline{\footnotesize bool, \code{true}} & bispectrum = true & Use interlaced grids to reduce aliasing \\[2pt]
\param{bispectrum_subtract_shotnoise}\newline{\footnotesize bool, \code{false}} & bispectrum = true & Subtract shot-noise \\[2pt]
\param{bispectrum_shell_overlap_weights}\newline{\footnotesize bool, \code{false}} & bispectrum = true & Weight the Fourier cells by their volume overlap with the k-shells (continuous shells, as in CONCEPT) instead of 0/1 from the cell centre \\[2pt]
\end{longtable}}

\section{Hybrid bias spectra}

{\small\begin{longtable}{@{}>{\raggedright\arraybackslash}p{4.3cm}>{\raggedright\arraybackslash}p{2.8cm}>{\raggedright\arraybackslash}p{8.1cm}@{}}
\toprule Parameter (type, default) & Used when & Description \\ \midrule
\endfirsthead
\toprule Parameter (type, default) & Used when & Description \\ \midrule
\endhead
\bottomrule
\endlastfoot
\param{hybridbias}\newline{\footnotesize bool, \code{false}} & always & Compute the hybrid bias (Lagrangian bias basis) spectra \code{P_ab}(k) at every output: the cross spectra of the operators 1, \code{delta_L}, \code{delta_L}\textasciicircum{}2, s\textasciicircum{}2, nabla\textasciicircum{}2 \code{delta_L} advected with the particles (Modi et al. 2020, Kokron et al. 2021). \code{delta_L} is the linear field of the IC scaled to the output redshift with the growth factor of the gravity model \\[2pt]
\param{hybridbias_nmesh}\newline{\footnotesize int, \emph{required}} & hybridbias = true & Grid size for the advected fields \\[2pt]
\param{hybridbias_output_linear_field}\newline{\footnotesize bool, \code{false}} & hybridbias = true & Write the linear field \code{delta_L}(\code{z_ini}) on the particle lattice to \code{hybridbias_linear_field_<name>.dat} in \code{output_folder} (FML linear field file, ComputePowerSpectra/HybridBiasLinearFieldIO.h): the input of the hybrid bias spectra of SnapshotAnalyzer (\code{hybridbias_linear_field_source} = "file") \\[2pt]
\param{hybridbias_rsd}\newline{\footnotesize bool, \code{false}} & hybridbias = true & Also the redshift-space multipoles \code{P_ab}\textasciicircum{}l(k) of the spectra (the operators assigned at the redshift-space positions of the particles, mean over the three axes): \code{hybridbias_rsd_z<z>.txt.} With \code{hybridbias_zeldovich_field} / \code{hybridbias_zeldovich} / \code{hybridbias_oneloop} also the Zel'dovich field (\code{hybridbias_rsd_zeldovichfield_z<z>.txt}), its exact mean (\code{hybridbias_rsd_zeldovich_z<z>.txt}), the control variate (\code{hybridbias_rsd_cv_z<z>.txt}) and the one-loop LPT prediction with EdS kernels (\code{hybridbias_rsd_oneloop_z<z>.txt}), with the growth rate f at k = 0.01 h/Mpc (\code{ComputePowerSpectra/HYBRID_BIAS_RSD.md}) \\[2pt]
\param{hybridbias_rsd_ellmax}\newline{\footnotesize int, \code{4}} & hybridbias = true & The largest multipole of the redshift-space spectra (even) \\[2pt]
\param{hybridbias_noperators}\newline{\footnotesize int, \code{5}} & hybridbias = true & Use the first noperators of (1, delta, delta\textasciicircum{}2, s\textasciicircum{}2, nabla\textasciicircum{}2 delta): 15 spectra for 5 \\[2pt]
\param{hybridbias_density_assignment_method}\newline{\footnotesize string, \code{"CIC"}} & hybridbias = true & Density assignment method. Options: \code{NGP}, \code{CIC}, \code{TSC}, \code{PCS}, \code{PQS} \\[2pt]
\param{hybridbias_interlacing}\newline{\footnotesize bool, \code{true}} & hybridbias = true & Use interlaced grids to reduce aliasing \\[2pt]
\param{hybridbias_subtract_shotnoise}\newline{\footnotesize bool, \code{false}} & hybridbias = true & Subtract the Poisson shot noise mean(\code{w_a} \code{w_b})/Npart (the usual basis spectra keep it) \\[2pt]
\param{hybridbias_smoothing_scale}\newline{\footnotesize double, \code{0}} & hybridbias = true & Smooth the linear field with this scale (Mpc/h) before forming the operators (0: none) \\[2pt]
\param{hybridbias_smoothing_filter}\newline{\footnotesize string, \code{"gaussian"}} & hybridbias = true & The smoothing filter. Options: \code{tophat}, \code{sharpk}, \code{gaussian} \\[2pt]
\param{hybridbias_kcut_over_nyquist}\newline{\footnotesize double, \code{1}} & hybridbias = true & Keep only the modes with k < this times the Nyquist frequency of the particle lattice in the linear field (a spherical cut, so the operators are isotropic as assumed in the LPT prediction). 0: keep all the modes of the lattice grid (a cube) \\[2pt]
\param{hybridbias_zeldovich}\newline{\footnotesize bool, \code{false}} & hybridbias = true & Also write the Zel'dovich (resummed 1LPT) prediction of the spectra with the same cut and smoothing of the operators: \code{hybridbias_zeldovich_z<z>.txt} \\[2pt]
\param{hybridbias_zeldovich_field}\newline{\footnotesize bool, \code{false}} & hybridbias = true & Also write the spectra of the Zel'dovich field of the same initial conditions (the lattice moved with the 1LPT displacement of the linear field at the output, with the same operator weights): \code{hybridbias_zeldovichfield_z<z>.txt.} With the Zel'dovich prediction (\code{hybridbias_zeldovich}, the mean of these spectra) this is the control variate of Kokron et al. 2022: \code{P_cv} = P - beta (\code{P_zeldovichfield} - \code{P_zeldovich}) has the same mean as P and much less sample variance on large scales \\[2pt]
\param{hybridbias_zeldovich_field_only}\newline{\footnotesize bool, \code{false}} & hybridbias = true & No time stepping: generate the initial conditions and write only the spectra of their Zel'dovich field at the output redshifts (as \code{hybridbias_zeldovich_field}). Cheap: for the mean of the control variate over many realizations in the same box, or for adding the control variate to existing runs (same seed and IC settings) \\[2pt]
\param{hybridbias_oneloop}\newline{\footnotesize bool, \code{false}} & hybridbias = true & Also write the one-loop LPT (CLEFT) prediction of the spectra (no counterterms): \code{hybridbias_oneloop_z<z>.txt} \\[2pt]
\param{hybridbias_lpt_kmin}\newline{\footnotesize double, \code{0}} & hybridbias = true & Lower limit (h/Mpc) of the k-integrals of the LPT predictions (0: the start of the linear P(k) table, or 1e-5). The box has no modes below its fundamental mode \\[2pt]
\param{hybridbias_zeldovich_cv_beta_k0}\newline{\footnotesize double, \code{0.618}} & hybridbias = true & The coefficient of the Zel'dovich control variate written to \code{hybridbias_cv_z<z>.txt} (with \code{hybridbias_zeldovich_field} and \code{hybridbias_zeldovich}): beta(k) = [1 - tanh((k - k0) / dk)] / 2 (Kokron et al. 2022, Eq. 4.7) with k0 (h/Mpc) this. 0: beta = 1 \\[2pt]
\param{hybridbias_zeldovich_cv_beta_dk}\newline{\footnotesize double, \code{0.167}} & hybridbias = true & The width dk (h/Mpc) of beta(k) of the Zel'dovich control variate (see \code{hybridbias_zeldovich_cv_beta_k0}) \\[2pt]
\param{hybridbias_oneloop_kernels}\newline{\footnotesize string, \code{"EdS"}} & hybridbias = true & The kernels of the one-loop prediction. EdS: the Einstein-de Sitter kernels with the linear P(k) of the gravity model. model: the scale- and time-dependent LPT kernels of the gravity model (f(R) and DGP with their screening if \code{gravity_model_screening}, the other models from \code{G_eff}(k,a) without screening, LCDM for GR). Options: \code{EdS}, \code{model} \\[2pt]
\end{longtable}}

\section{Images}

{\small\begin{longtable}{@{}>{\raggedright\arraybackslash}p{4.3cm}>{\raggedright\arraybackslash}p{2.8cm}>{\raggedright\arraybackslash}p{8.1cm}@{}}
\toprule Parameter (type, default) & Used when & Description \\ \midrule
\endfirsthead
\toprule Parameter (type, default) & Used when & Description \\ \midrule
\endhead
\bottomrule
\endlastfoot
\param{images}\newline{\footnotesize bool, \code{false}} & always & Draw PNG images of fields at every output: a slab projected along an axis (\code{image_<field>_z<z>.png} in the snapshot folder) \\[2pt]
\param{images_fields}\newline{\footnotesize string list, \texttt{\{density\}}} & images = true & The fields to draw, e.g. \{"density", "\code{velocity_divergence}"\}: density (1 + delta of the particles, log colours), \code{velocity_divergence} (-div v / (a H) = f delta in linear theory); SimulationCode also neutrinos (\code{delta_nu}), \code{baryons_minus_cdm} (\code{delta_b} - \code{delta_c}) and \code{fifth_force} (|\code{F_5th}| / |\code{F_Newton}|) \\[2pt]
\param{images_nmesh}\newline{\footnotesize int, \code{256}} & images = true & Grid size of the fields \\[2pt]
\param{images_npix}\newline{\footnotesize int, \code{0}} & images = true & Pixels per side of the images (0: nmesh) \\[2pt]
\param{images_axis}\newline{\footnotesize int, \code{2}} & images = true & Projected along this axis (0, 1, 2: x, y, z) \\[2pt]
\param{images_slab_center}\newline{\footnotesize double, \code{0.5}} & images = true & Centre of the projected slab along the axis (units of the box) \\[2pt]
\param{images_slab_thickness}\newline{\footnotesize double, \code{0.1}} & images = true & Thickness of the projected slab (units of the box; 1: the whole box). The image is the mean over it \\[2pt]
\param{images_density_assignment_method}\newline{\footnotesize string, \code{"CIC"}} & images = true & Density assignment method (NGP, CIC, TSC, PCS, PQS) \\[2pt]
\param{images_colormap}\newline{\footnotesize string, \code{""}} & images = true & Colour map for all fields (magma, inferno, viridis, plasma, cividis, coolwarm, grey, cubehelix; \code{_r} reverses). Empty: each field its own (density magma, signed fields coolwarm) \\[2pt]
\param{images_npy}\newline{\footnotesize bool, \code{false}} & images = true & Also write the projected values as .npy files (for numpy) \\[2pt]
\end{longtable}}

\section{Options with separate particle species}

Not everything is implemented yet for CDM and baryons as separate particles (\code{particle_baryons}) and for neutrino particles (\code{particle_neutrinos}). The code stops with a message if one of these is not satisfied (generated from \code{species_requirements()} in \code{src/SimulationParameters.h}):

{\small\begin{longtable}{@{}>{\raggedright\arraybackslash}p{9.0cm}cc@{}}
\toprule Needs & \code{particle_baryons} & \code{particle_neutrinos} \\ \midrule
\endfirsthead
\toprule Needs & \code{particle_baryons} & \code{particle_neutrinos} \\ \midrule
\endhead
\bottomrule
\endlastfoot
\code{timestep_method = Quinn} (the relativistic factors of each neutrino are applied on the plain kick and drift; FastPM, Tassev and BullFrog take their factors from the growth of a cold fluid, which the free-streaming neutrinos do not follow) &  & yes \\
\code{simulation_use_cola = false} (the LPT frame of the CDM+baryons does not describe the neutrinos: their thermal motion is not small relative to it) &  & yes \\
\code{timestep_method} other than \code{Tassev} with COLA (the modified Tassev operators get the baryon-CDM relative velocity wrong: +30\% in \code{P_b}-c with 20 linear steps) & yes &  \\
\code{simulation_cola_LPT_order} <= \code{ic_LPT_order} (with COLA) & yes &  \\
\code{simulation_enforce_LPT_trajectories = false} & yes & yes \\
\code{ic_type_of_input = camb}, \code{class}, \code{internal} or \code{tables} (species transfers at zini) & yes & yes \\
\code{ic_random_field_type = gaussian} or \code{nongaussian} & yes & yes \\
\code{ic_LPT_order = 1} or \code{2} & yes &  \\
\code{force_linear_massive_neutrinos = false} (the particles replace the linear neutrinos) &  & yes \\
\code{force_density_method = assignment} or \code{nufft} (not the sheet densities) & yes & yes \\
\code{force_nbody_gauge = false} & yes & yes \\
\code{lightcone = false} & yes & yes \\
\code{pofk_multipole = false} & yes & yes \\
\code{bispectrum = false} & yes & yes \\
\code{hybridbias = false} & yes &  \\
\end{longtable}}

